\documentclass[11pt,letterpaper]{article}
\PassOptionsToPackage{hyphens}{url}
\usepackage[margin=.83in,top=.8in,bottom=.8in,headheight=15pt]{geometry}
\usepackage[T1]{fontenc}
\usepackage{lmodern,amsmath,amssymb,mathtools,bm,booktabs,array,graphicx,tikz}
\usetikzlibrary{arrows.meta,positioning,calc,decorations.pathreplacing,patterns}
\definecolor{ink}{HTML}{17344D}\definecolor{accent}{HTML}{087E8B}\definecolor{pale}{HTML}{EDF5F5}\definecolor{gold}{HTML}{AD6B19}
\providecommand{\dd}{\,\mathrm{d}}\providecommand{\R}{\mathbb{R}}
\providecommand{\E}{\mathbb{E}}\providecommand{\Var}{\operatorname{Var}}
\providecommand{\Cov}{\operatorname{Cov}}\providecommand{\diag}{\operatorname{diag}}
\providecommand{\norm}[1]{\left\lVert#1\right\rVert}
\providecommand{\abs}[1]{\left|#1\right|}
\usepackage{microtype,hyperref,fancyhdr,enumitem,amsthm,sectsty,needspace}
\allsectionsfont{\raggedright}
\hypersetup{colorlinks=true,linkcolor=ink,urlcolor=accent,pdfauthor={Zan Ahmad},pdftitle={Surrogate error, amortized cost and training-step bias / Review}}
\setlength{\parindent}{0pt}\setlength{\parskip}{6pt}\linespread{1.07}
\setlist{itemsep=3pt,topsep=4pt,leftmargin=*}
\emergencystretch=2em\clubpenalty=10000\widowpenalty=10000
\pagestyle{fancy}\fancyhf{}\fancyhead[L]{\small Computational Cardiology}\fancyhead[R]{\small Lecture 19}
\fancyfoot[L]{\footnotesize Zan Ahmad}\fancyfoot[R]{\footnotesize\thepage}\renewcommand{\headrulewidth}{.3pt}
\newtheorem{proposition}{Proposition}\newtheorem{theorem}{Theorem}\newtheorem{lemma}{Lemma}
\theoremstyle{definition}\newtheorem{definition}{Definition}\newtheorem{example}{Example}
\numberwithin{equation}{section}
\begin{document}
\section*{Surrogate error, amortized cost and training-step bias --- review sheet}
Use this after attempting the worked examples and practice. Explain each claim in words before calculating.

\section*{Key ideas}
\begin{itemize}
\item A useful claim includes regime, observable and acceptance threshold.
\item Setup and per-query cost answer different questions.
\item Reference accuracy must be checked independently.
\item A learned generator can depend on the training discretization.
\item A computational benchmark is not biological validation.
\end{itemize}
\section*{Common errors and repairs}
\textbf{Reporting only a favorable warm inference time.} State the boundary and separately account for setup and geometry costs.

\textbf{Assuming a higher-order integrator always improves target accuracy.} Test the learned equation and training/inference discretizations together.

\textbf{Treating reference labels as error-free reality.} Separate reference verification from model validation.

\section*{Retrieve without your notes}
\begin{itemize}
\item When does a positive setup break-even threshold exist?
\item What can a smaller step converge toward?
\end{itemize}
\section*{Readiness check}
\begin{itemize}
\item Separate surrogate error, reference-solver error and biological validation evidence.
\item Measure query costs with an explicit timing boundary and compute a setup break-even condition.
\item Explain how a fitted generator can depend on its training time step.
\item Design a falsifiable, bounded comparison with grouped test regimes and failure criteria.
\end{itemize}
\end{document}
